% Owned implication fragment for C.3bt.  The marked edge-family lemma is
% proved in the preceding cycles fragment.  No independent result row.
\subparagraph{From marked edge families to the whole component kernel.}
We use the notation and exact exceptional threefold scope of C.3bm and
C.3bo.  In particular $E=\bigcup_iE_i$ is a projective SNC surface over
an algebraically closed field of characteristic $p>0$, and a positive
exceptional combination $D=\sum_i a_iE_i$ satisfies
$\mathcal O_E(-D)$ ample.  Put
\[
 S=F_\gamma^2K_0(E)_{\mathbf Q},\qquad J=J_K(E),\qquad
 Q=\theta\operatorname{CH}_1(E)_{\mathbf Q},\qquad
 \varphi(s)=\bigl(\operatorname{ch}_2(s|_{E_i})\bigr)_i.
\]
Here $J$ is generated by the endpoint differences of actual perfect
gamma-two classes on $E$ times smooth connected projective curves, as
in C.3bm.  We prove, using the marked edge-family lemma above, that
\begin{equation}
 \varphi:S\twoheadrightarrow\bigoplus_i\operatorname{CH}_0(E_i)_{\mathbf Q},
 \qquad \ker\varphi\subset J.                         \tag{UK.1}
\end{equation}
Neither a dual-complex cohomology hypothesis nor a triple-point period
hypothesis is imposed.

\emph{1. Keep the actual totalization and prove pointwise surjectivity.}
For a fixed smooth projective curve $P$, or $P=\operatorname{Spec}k$,
let
\[
 \mathcal C_P^0=\bigoplus_iK(E_i\times P)_{\mathbf Q},\quad
 \mathcal C_P^1=\bigoplus_eK(C_e\times P)_{\mathbf Q},\quad
 \mathcal C_P^2=\bigoplus_vK(P)_{\mathbf Q},\quad
 \mathcal T_P=\operatorname{fib}(\mathcal C_P^1\to\mathcal C_P^2).
\]
The canonical geometric nullhomotopy, not a replacement chosen on
homotopy groups, gives
\begin{equation}
 K(E\times P)_{\mathbf Q}\simeq
          \operatorname{fib}(\mathcal C_P^0\to\mathcal T_P).
                                                               \tag{UK.2}
\end{equation}
This is the every-degree closed-union descent proved in (ET.6), from
\citet[Theorems A and C, printed pp.~29--30]{rootGeisserHesselholt2011}
and \citet[Theorem 8.1 and proof, p.~367]{D_ThomasonTrobaugh1990}.
The exact sequences, natural under both endpoint sections, are
\begin{align}
 \bigoplus_vK_2(P)_{\mathbf Q}&\longrightarrow\pi_1\mathcal T_P
 \longrightarrow\bigoplus_eK_1(C_e\times P)_{\mathbf Q}
 \longrightarrow\bigoplus_vK_1(P)_{\mathbf Q},             \tag{UK.3}\\
 \pi_1\mathcal T_P&\xrightarrow{\partial_P}K_0(E\times P)_{\mathbf Q}
 \longrightarrow\bigoplus_iK_0(E_i\times P)_{\mathbf Q}.
                                                               \tag{UK.4}
\end{align}

Every tuple of component zero-cycles has a unique pure weight-two class
\[
                         b\in\bigoplus_iK_0(E_i)_{\mathbf Q}.
\]
Its obstruction
$a\in\pi_0\mathcal T_k$ has eigenvalue $4$ for $\psi^2$ and maps to
zero in $\bigoplus_eK_0(C_e)_{\mathbf Q}$, whose weights are $0$ and $1$.
The first fibre sequence in (UK.2) therefore puts $a$ in the image of
$\bigoplus_vK_1(k)_{\mathbf Q}$, on which $\psi^2$ is multiplication by
$2$.  Consequently $4a=2a$, and $a=0$.  This uses only naturality of
the operations and exactness, not a semisimple weight decomposition of
$\pi_0\mathcal T_k$.  Equivalently that group is annihilated by
$(\psi^2-1)(\psi^2-2)^2$, which is invertible at the eigenvalue $4$.
Exactness lifts $b$ to $K_0(E)_{\mathbf Q}$; the rational weight-two
projector places the lift in $S$.  The finite gamma filtration and the
projectors on the projective surface and its products are supplied by
\citet[Chapter II, Proposition 4.9, Theorem 4.10, and Remark 8.8.5]
{D_Weibel2013KbookII}.  This proves the surjectivity in (UK.1).

\emph{2. Reduce every component-kernel class to degree-zero edge terms.}
Take $s\in\ker\varphi$.  The rational Chern character on the smooth
components makes its actual component $K_0$ restrictions zero, so
(UK.4) gives $y\in\pi_1\mathcal T_k$ with $\partial_k y=s$.
Replace $y$ by $(\psi^2-2)y/2$.  Its boundary remains $s$; on each
smooth proper double curve its edge image has weight two and is a
finite sum of terms $u([L]-1)$, with $u\in k^*$ and
$L\in\operatorname{Pic}(C_e)$.
This is the localization and Picard-times-units generation used in
C.3br and C.3bq
\citep[Section 7, Propositions 2.10 and 3.2]{rootQuillen1973}
\citep[Lemma 5.2 and its full proof, pp.~395--396]{D_PedriniWeibel1994}.
Commutation with the singular closed-union boundary uses the
arbitrary-scheme iterated-fibre and delooping construction of
\citet[p.~258, paragraph preceding Lemma 2.1]{frC3bcLevine1994};
we do not apply that lemma's regular-ambient hypothesis to $E$.

The weighted identity in (ET.9) makes alternating restriction onto
the edge degrees:
\[
 \bigoplus_i\operatorname{Pic}(E_i)_{\mathbf Q}
                  \twoheadrightarrow\bigoplus_e\mathbf Q.
\]
For each of the finitely many units, multiply component
divisor lifts of the edge degree vector by that constant unit in
$K_1$.  Subtract their canonical images in $\pi_1\mathcal T_k$.
Their boundaries in (UK.4) are zero.  The remaining edge image of $y$
is a finite sum of vectors
\begin{equation}
       \gamma_e=u([L_e]-1),\qquad
       L_e\in\operatorname{Pic}^0(C_e)(k),                \tag{UK.5}
\end{equation}
one unit $u$ at a time.  Clearing rational denominators before the
construction and dividing afterwards is harmless.  No surjectivity
of the full Picard restriction map has been assumed.

\emph{3. The marked families remove the remaining edge vectors.}
Apply the preceding marked edge-family lemma to (UK.5), with every
triple point on each connected edge included among the markings.
It supplies a smooth connected projective curve $P$ with endpoints
$o,t$ and actual classes
$\Gamma_e\in K_1(C_e\times P)_{\mathbf Q}$ such that
\begin{equation}
 (\Gamma_e)_t-(\Gamma_e)_o=\gamma_e,\qquad
 \Gamma_e|_{q_{e,v}\times P}=0\quad\hbox{in }K_1(P)_{\mathbf Q}
                  \quad\hbox{for every }v\supset e.       \tag{UK.6}
\end{equation}
The second equality is an equality on the actual proper parameter,
not just equality at its endpoints.  Thus (UK.3) gives
$\widetilde y\in\pi_1\mathcal T_P$ with edge image
$(\Gamma_e)_e$.  The class $\partial_P\widetilde y$ is an actual
perfect $K_0$ family.  On the three-dimensional projective scheme
$E\times P$, apply the rational Adams-weight-two projector to it.
Its endpoint difference therefore lies in $J$ and has the same
weight-two component-kernel contribution as
$\partial_k(\widetilde y_t-\widetilde y_o)$.

The difference between the originally chosen $y$ and these endpoint
lifts has zero edge image.  Exactness of (UK.3), with $P=k$, places
it in the image of $\bigoplus_vK_2(k)_{\mathbf Q}$.  The canonical
double boundary of this image is contained in $J$: the actual
proper-curve symbol construction in C.3bq Step~2 realizes every
Milnor symbol as a difference of evaluations of a class in $K_2$ of
a smooth proper curve, and naturality of (UK.3)--(UK.4) produces the
required family.  Its weight-two projection retains the desired
endpoint class.  This treats the full coherent-lift ambiguity;
no unique lift in (UK.3) is asserted or required.  Finite sums of
(UK.5) now prove $s\in J$, establishing (UK.1).

\emph{4. The precise component Albanese-kernel quotient.}
For each component write
\[
 T(E_i)=\ker\bigl(\operatorname{CH}_0(E_i)^0_{\mathbf Q}
       \longrightarrow\operatorname{Alb}(E_i)(k)\otimes\mathbf Q\bigr),
 \qquad T_E=\bigoplus_iT(E_i).
\]
Let $S_{00}$ have the meaning of C.3bo: both all component degrees
and all component Albanese images vanish.  Surjectivity in (UK.1)
gives $\varphi(S_{00})=T_E$.  Its kernel is contained in
$J\cap S_{00}$.  The exact denominator splitting of (AM.4) therefore
gives
\begin{equation}
 B^2(Z\setminus E)_{\mathbf Q}\simeq
 \frac{T_E}{
   \varphi(J\cap S_{00})+\varphi(Q\cap S_{00})}.          \tag{UK.7}
\end{equation}
The two denominator terms retain the induced actual proper-perfect
families and the actual exceptional cycle image, respectively.
For completeness, (AM.4) uses the Poincare subgroup
$U\subset J\cap Q\cap S_0$ whose component Albanese map is onto:
subtracting $u\in U$ first removes the Albanese coordinate, and for
$j+q\in S_{00}$ the replacement
$j+q=(j+u)+(q-u)$, with $\alpha_E(u)=\alpha_E(q)$, separates the
two denominator summands inside $S_{00}$.  Thus (UK.7) is not an
informal quotient by a sum of separately projected coordinates.

In particular, if all $T(E_i)$ vanish, the ordinary rational quotient
vanishes with no condition on $H^2(\Delta,\mathbf Q)$ or on triple
periods.  For arbitrary components the remaining group is exactly
(UK.7).  No assertion that all component Albanese kernels are induced
by globally compatible proper perfect families is made.
